% CECI EST UN FICHIER DE MODÈLE LATEX POUR LES ARTICLES INCLUS DANS
% *Anthology of Computers and the Humanities*.
% 'final' LORS DE LA CRÉATION DE LA VERSION FINALE DE L'ARTICLE.
% NE PAS modifier la classe du document
%\documentclass[final,fra]{anthology-ch} % pour la version finale
\documentclass[final,fra]{anthology-ch}         % pour la soumission

% CHARGER LES PACKAGES LaTeX
\usepackage{booktabs}
\usepackage{graphicx}

% AJOUTEZ vos propres packages avec \usepackage{}

% TITRE DE LA SOUMISSION
% Remplacez ceci par le titre de votre soumission
\title{Humanités numériques et justice épistémique africaine : le cas du transfert culturel de l’art musical Bambala dans la musique congolaise moderne}

% INFORMATIONS SUR LES AUTEURS ET LES AFFILIATIONS
% Pour chaque auteur, utilisez une nouvelle commande \author, avec
% les numéros entre crochets indiquant les affiliations correspondantes
% (section suivante) et l'ORCID-ID pour chaque auteur.
\author[1]{Kageri-Andy Kindongo}[
  orcid=0009-0003-8762-5715,
  email=kindongok@edu.univ-eiffel.fr
]

\author[2]{Saidou Kodio}[
  orcid=0009-0004-5899-2708,
  email=saidou.kodio@u-pec.fr
]

% Il doit y avoir un appel à \affiliation pour chaque affiliation des auteurs.
% Plusieurs affiliations peuvent être attribuées à un auteur
% et une affiliation peut être attribuée à plusieurs auteurs.
%\affiliation{1}{Université Gustave Eiffel LISAA - Laboratoire Littératures, Savoirs et Arts, École doctorale cultures et sociétés, Champs-Sur-Marne, France}
\affiliation{1}{Université Gustave Eiffel / LISAA (EA 4120), Champs-Sur-Marne, France}
%\affiliation{2}{Université Paris-Est Créteil Val-de-Marne - Paris 12 École doctorale cultures et sociétés, LIS, Créteil, France}
\affiliation{2}{Université Paris-Est Créteil Val-de-Marne - Paris 12 / LIS (EA 4395), Créteil, France}

% MOTS-CLÉS
% Fournissez un ou plusieurs mots-clés séparés par des virgules
% en utilisant les commandes suivantes
\keywords[fra]{humanités numériques, justice épistémique, ethnomusicologie, transfert culturel, Bambala, musique congolaise, études postcoloniales}
\keywords[eng]{digital humanities,
epistemic justice, ethnomusicology, cultural transfer; Bambala, Congolese music, postcolonial studies}

% MÉTADONNÉES DE LA PUBLICATION
% Ces champs seront remplis lors de la publication ;
% ils peuvent être laissés par défaut lors de la soumission
\pubyear{2026}
\pubvolume{4}
\pagestart{227}
\pageend{234}
\paperorder{28}
\conferencename{Actes de la Conférence Humanistica}
\conferenceeditors{Serena Crespi, Simon Gabay, Martin Grandjean, Ariane Pinche, Marie Puren et Léa Saint-Raymond}
\doi{10.63744/wO4Pvylm5glD}
\customhead{}

\addbibresource{bibliography.bib}

%%%%%%%%%%%%%%%%%%%%%%%%%%%%%%%%%%%%%%%%%%%%%%%%%%%%%%%%%%%%%%%%%%%%%%%%%%%
% DÉBUT DU TEXTE
\begin{document}

\maketitle

\begin{abstract}
    This article examines how Digital Humanities can contribute to restoring epistemic justice toward African artistic knowledge, which has long been marginalized in global intellectual history. Focusing on the musical traditions of the Bambala people (Democratic Republic of Congo), whose rhythms and dances have significantly influenced modern Congolese music without explicit recognition, we analyze the mechanisms of cultural transfer and the forms of invisibilities that accompany them. By articulating postcolonial philosophy, ethnomusicology, and Digital Humanities, we propose a theoretical and methodological framework for documenting, contextualizing, and digitally preserving these musical practices. The study highlights the need for ethical, legal, and technical infrastructures that ensure the visibility and legitimacy of local knowledge within global digital ecosystems.
\end{abstract}

\section{Introduction}\label{introduction-1}

L'histoire intellectuelle de l'Afrique subsaharienne a été marquée par une marginalisation persistante de ses savoirs, souvent interprétés à travers des catégories hiérarchisantes. La prédominance de l'oralité a été assimilée à une absence d'abstraction, conduisant certains penseurs européens à reléguer les sociétés africaines hors du champ de l'histoire universelle \cite{hegel2018phenomenologie} ou à qualifier leurs modes de pensée de \og prélogiques\fg{} \cite{levybruhl1951fonctions}. Ces lectures ont contribué à ce que Miranda Fricker \cite{fricker2007epistemic} nomme une injustice épistémique, c'est-à-dire une disqualification structurelle des savoirs d'un groupe social, soit par préjugés (injustice testimoniale), soit par absence de cadres conceptuels permettant leur reconnaissance (injustice herméneutique).

Dans ce contexte, les humanités numériques constituent un espace critique où se redéfinissent les modalités de production, de diffusion et de légitimation du savoir. Elles offrent des outils pour documenter, archiver et rendre visibles des formes de connaissances longtemps invisibilisées, tout en favorisant leur mise en dialogue avec d'autres traditions intellectuelles. Cette perspective rejoint la philosophie de la traduction et de l'hospitalité des savoirs défendue par Souleymane Bachir Diagne \cite{diagne2025universels}, qui appelle à dépasser l'\og universalisme impérial pour construire une universalité plurielle.\fg{}

Le cas de l'art musical des Bambala de la République Démocratique du Congo illustre de manière exemplaire ces enjeux. Bien que leurs pratiques musicales aient profondément influencé la musique congolaise moderne, ces transferts culturels demeurent rarement reconnus ou contextualisés. L'absence de documentation, de référencement et de visibilité numérique contribue à une forme d'injustice épistémique, où les savoirs mbala circulent sans que leurs détenteurs soient identifiés comme sources légitimes.

Comment les humanités numériques peuvent-elles contribuer à la restauration de la justice épistémique autour de l'art musical mbala, dont les transferts culturels vers la musique congolaise moderne restent largement invisibilisés? Cet article propose d'examiner comment les humanités numériques peuvent contribuer à restaurer cette justice épistémique. Nous présentons d'abord un cadre théorique articulant postcolonialité, justice épistémique et humanités numériques. Nous analysons ensuite les caractéristiques de l'art musical mbala et les mécanismes de son transfert vers la musique congolaise moderne. Enfin, nous ouvrons des perspectives pour une patrimonialisation numérique durable et inclusive.

\section{Humanités numériques et justice épistémique : un cadre théorique}

\subsection{Hiérarchies épistémiques et marginalisation des savoirs africains}

La pensée hégélienne, en excluant l'Afrique subsaharienne du mouvement
de l'esprit vers la liberté, a contribué à installer une hiérarchie épistémologique durable \cite{hegel2018phenomenologie}. L.~Lévy-Bruhl \cite{levybruhl1951fonctions}, en qualifiant certaines sociétés de \og prélogiques\fg{}, a prolongé cette vision en opposant rationalité occidentale et pensée \og mystique\fg{}. Ces approches ont nourri ce que Diagne \cite{diagne2025universels} nomme \og  l'universalisme impérial\fg{}, c'est-à-dire l'imposition des catégories européennes comme normes universelles. De plus, Levi-Strauss \cite{leviStrauss_pensee_1962} défend le fait  que chaque culture constitue un système de connaissance à part entière. Dès lors, l’absence de documentation, d’études ou de transmission de ces patrimoines culturels contribuerait à invisibiliser des formes légitimes de connaissance, pouvant être interprétée comme une forme d’injustice épistémique.

La notion d’injustice épistémique proposée par Fricker \cite{fricker2007epistemic} permet d’éclairer ces mécanismes. L’injustice testimoniale survient lorsque des préjugés d’identité diminuent la crédibilité d’un locuteur. L’injustice herméneutique apparaît lorsqu’un groupe social ne contribue pas à la production des concepts partagés, ce qui empêche la reconnaissance de ses expériences. Ces deux formes d’injustice se retrouvent dans la réception des savoirs artistiques africains, souvent décontextualisés, anonymisés ou attribués à des catégories génériques (\og tradition\fg{}, \og folklore\fg{}).

\subsection{Hiérarchies épistémiques et marginalisation des savoirs africains}

Les humanités numériques ne se limitent pas à l'usage d'outils technologiques : elles constituent un champ critique où se redéfinissent les critères de production et de validation du savoir. Elles peuvent contribuer à la justice épistémique en rendant visibles des savoirs marginalisés, en contextualisant les œuvres, en associant les communautés d’origine, et en favorisant la circulation plurielle des connaissances. Cependant, elles peuvent également reconduire les hiérarchies existantes si les corpus numérisés ne sont pas analysés dans leur contexte culturel et symbolique. La numérisation doit donc s’accompagner d’une réflexion éthique, juridique et politique.

\subsection{Application au cas africain : vers une épistémologie plurielle et traductive}

Une épistémologie plurielle et traductive, selon Mboa Nkoudou \cite{nkoudou2015strategies}, consiste à reconnaître la diversité des savoirs africains et à dépasser les cadres hérités du colonialisme: chez les Mbala, la musique, les rites et la parole sont des formes de connaissance qui organisent le sens et la cohésion sociale \cite{matangila1994rites}, dans un contexte historique contemporain marqué par les transformations sociales, politiques et culturelles décrites par Kiangu Sindani \cite{kiangu_sindani_kwilu_2009}. Cette approche permet de contribuer à une justice épistémique, en rendant visibles et intelligibles les savoirs mbala appropriés, réinterprétés et diffusés sans reconnaissance de leurs producteurs.

\section{L'art musical mbala: pratiques, instruments et fonctions sociales}

\subsection{Présentation des Bambala et méthodologie}

Les Bambala, peuple du Kwango et du Kwilu (RDC), parlent des variantes du Gimbala/Kimbala, le Kikongo ya leta et le français. Notre étude se concentre sur ceux du territoire de Masi-Manimba, avec une méthodologie relevant de l'anthropologie de la musique, fondée sur l'écoute active, la collecte de données, l'analyse, l'archivage et la valorisation scientifique.

\subsection{Caractéristiques de la musique mbala}

Torday et Joyce \cite{torday1905notes} décrivent la musique mbala comme profondément intégrée à la vie sociale et rituelle. Le Malunga, par exemple, est exécuté lors des rites funéraires.

\textbf{Voix et structures vocales}

Les instruments principaux sont:

\begin{itemize}
    \item le Mungungi, tambour central de la communauté \cite{mudindaambi1976objets}
    \item le Kapetengi, tambour accompagnateur également cylindrique selon Lumbwe Mudindaambi \cite{mudindaambi1976objets}.
    \item le Gikedi, cloche rituelle \cite{mudindaambi1976objets}.
    \item des sifflets de chasse ou de guerre. Ces instruments structurent les danses et les chants, notamment les Lugamba, Kidagudagu, Bwandji ou Mapotopoto.
\end{itemize}

\section{Transfert culturel vers la musique congolaise moderne}

\subsection{Présentation des Bambala et méthodologie}

Pour Malinowski \cite{malinowski1945dynamics}, le transfert culturel adopte des éléments d'une autre culture quand des groupes sont en contact, transformant peu à peu la culture d'origine. Notre analyse est basée sur la deuxième école de l’anthropologie de la musique  \cite{ayats_tout_1996} et distingue deux modes de transfert:

\begin{itemize}
    \item L'imitation, où les motifs musicaux sont repris sans mention de leur origine.
    \item La référence, où l'œuvre source est citée explicitement.
\end{itemize}

\noindent Ces transferts s'accompagnent souvent d'une décontextualisation: des rythmes rituels deviennent des éléments de divertissement.

\subsection{Œuvres sources et influences}

Nous avons identifié cinq œuvres mbala ayant inspiré vingt chansons congolaises modernes, consultables sur les plateformes numériques. Ces influences concernent notamment:

\begin{itemize}
    \item \textit{Kidagudagu} \cite{mbanga_mumbungu_gi_1982};
    \item \textit{Makwa Ndumbu} \cite{mavula_bankwa_1957};
    \item \textit{Mumembu} \cite{mapombu_mumembu_nodate};
    \item \textit{Souvenirs ya indépendance} \cite{mavula_souvenir_2006};
    \item \textit{Kasongo Muwaya} \cite{mbanga_mumbungu_kasongo_1986},
\end{itemize}

Ces œuvres ont été reprises par des artistes tels que ceux des orchestres Wenge Musica et Grand Zaïko, Mbilia Bel, Popaul Amisi, Bopol Mansiamina, Tshanda Sangwa (Bb Rico), Flamme Kapaya ou Suzuki Luzubu. \subsection{Typologie des transformations musicales: du rituel mbala à la musique populaire congolaise}

La majorité de chansons congolaises modernes a fait l'objet d'un transfert par imitation, contre 6 par référence.

\begin{table}[htp]
    \centering\small
    \begin{tabular}{p{\dimexpr\linewidth/3 - 2\tabcolsep\relax}p{\dimexpr\linewidth/3 - 2\tabcolsep\relax}p{\dimexpr\linewidth/3 - 2\tabcolsep\relax}}
    \toprule
    Nom Titre Original
     & Œuvre Originale
     & Œuvres Modernes
     \\
    \midrule
    \emph{Kidagudagu} (1982) & Danse Kidagudagu.
    Ré Majeur, tempo 125. & Même danse, Tempo moins rapide.
    Tonalité majeure, pas de référence \\
    \emph{Makwa Ndumbu} (1957) & Danse Makwa Ndumbu.
    Mib Majeur, tempo 125.
    
    Mesure 2/4 & Même danse, Tempo moins rapide.
    
    Do Majeur, tempo 108
    
    Mesure 4/4, références titre et auteur.
    \\
    \emph{Mumembu} (Mapombu) & Danse Mapotopoto
    
    Tempo rapide & Même danse.
    
    Tempo rapide, référence titre.
    \\
    \emph{Souvenirs ya indépendance} (1957) & Danse Bwandji, mesure 4/4
    
    Do Majeur,Tempo 115. & Mêmes danse, Tempo et tonalité (Do M).
    Références titre et auteur. \\
    \emph{Kasongo Muwaya} (1986) & Kidagudagu, Sib Majeur.
    Mesure 4/4, tempo 125. & Même danse et mesure, tempo moins rapide.
    
    Sol Majeur.
    Aucune référence \\
    \bottomrule
    \end{tabular}
    \caption{Tableau des transferts culturels observés.}
    \label{tab:tab1}
    \vspace{-.3cm}
\end{table}

\subsection{Injustice épistémique dans la circulation musicale}

\begin{table}[htp]
    \centering\small
    \begin{tabular}{p{\dimexpr\linewidth/3 - 2\tabcolsep\relax}p{\dimexpr\linewidth/3 - 2\tabcolsep\relax}p{\dimexpr\linewidth/3 - 2\tabcolsep\relax}}
    \toprule
    Type d'injustice
     & Manifestation chez les Bambala
     & Solutions numériques possibles
     \\
    \midrule
    \emph{Injustice testimoniale} & Œuvres originales reprises mais non
    reconnues, noms effacés.
    Artistes modernes: Si œuvre en Gimbala, disqualifiée car \og trop
    traditionnelle\fg{}.
    
    Versions partiellement en Gimbala = seules entendues à grande échelle.
    Auditeurs : interprétation culturelle ignorée. & Archives numériques
    communautaires (plateformes où les Mbala déposent, datent et signent
    leurs rythmes, récits et performances).
    Plateformes curatoriales alternatives.
    
    Bases de données de traçabilité (origines visibles).
    
    Archives et espaces d'écoute.
    \\
    \midrule
    \emph{Injustice herméneutique} & Artistes originaux : peu d'accès aux
    cadres juridiques, médiatiques.
    Artistes modernes marginalisés : ressent l'injustice mais ne peut pas la
    faire reconnaître.
    Artistes modernes : ne voient pas l'injustice qu'ils reproduisent
    
    Auditeurs : manière de comprendre la musique marginalisée. & Ateliers,
    récits numériques participatifs.
    Cartographies interactives, visualisations des flux culturels
    
    Revues, podcasts et archives critiques numériques.
    
    Dispositifs de mémoire et de narration numérique.
    \\
    \bottomrule
    \end{tabular}
    \caption{Tableau des injustices épistémiques observés chez les Bambala.}
    \label{tab:tab2}
    \vspace{-.3cm}
\end{table}

Nous observons que le faible transfert par référence contribue à l'invisibilisation des sources, engendrant une injustice épistémique. L'absence de citation, de droits, de contextualisation, contribue ainsi à l'effacement des communautés d'origine.

Les humanités numériques rendent visibles les origines et les voix des cultures minorisées grâce aux archives et à la traçabilité. Elles permettent aussi de nommer et comprendre les injustices liées aux transferts culturels.

\section{Perspectives : vers une patrimonialisation numérique du savoir des Bambala}

\subsection{Approches interdisciplinaires}

L'étude du savoir des Bambala nécessite des méthodes combinées :

\begin{itemize}
    \item Des analyses qualitatives issues de la musicologie et de l'ethnomusicologie \cite{laliberte2009prealables,anakesa2016occident};
    \item Des analyses quantitatives appuyées sur corpus et patterns des musiques populaires \cite{panourgia_field_2021};
    \item Une mise en perspective historique des circulations culturelles \cite{bouchet_introduction_2016};
    \item Une patrimonialisation durable, qui suppose : la numérisation des corpus \cite{baudorre_mauriac_2020}, avec une structure de métadonnées \cite{renault_harmoniser_2020}, permettant ainsi leur valorisation par différents usages tels que ceux de la recherche \cite{michaud_linguistes_2024}, comme le proposent Carlin et Laborderie \cite{carlin2021bnf}, leur catalogage et la création de bases de données multimédias \cite{witz2025interoperabilite}, l’annotation des archives sonores \cite{george2025avaa} et la participation des communautés mbala \cite{panourgia_field_2021}.
\end{itemize}

\noindent Elle gagne enfin à intégrer une réflexion sur l'écologie du numérique, afin d'assurer la durabilité et la responsabilité des dispositifs d'archivage et de diffusion.

\subsection{Application au corpus mbala}

Appliquée à un corpus de plus de 150 enregistrements sonores des folklores des Bambala, en cours de structuration et de cartographie, cette approche contribue à l'identification des circulations et transformations des rythmes traditionnels dans les musiques populaires urbaines.
    \begin{figure}[h]
        \centering
        \includegraphics[width=0.9\textwidth]{./media/image1.png}
        \caption{Illustration de la cartographie des transferts culturels du folklore des Bambala~\cite{kindongo_cartographie_2026}.}
    \end{figure}

Elle permet de croiser analyses musicologiques qualitatives et analyses de patterns, en reliant les formes musicales aux contextes sociaux et  rituels de leur production.

\subsection{Patrimonialisation communautaire~: le cas des Archives Mbala}

Parallèlement, une démarche de patrimonialisation communautaire est menée à travers l'association de loi 1901 «~Archives Mbala~», menant des actions de conservation et de valorisation du patrimoine culturel des Bambala: ils ont déjà collecté plus de 500 enregistrements audiovisuels, ouvrages documentaires. La numérisation, l'annotation et la documentation de ces archives visent à restituer les savoirs musicaux dans leur contexte culturel et à renforcer la visibilité des acteurs mbala.
Ces actions ont permis de retrouver l'artiste Ya Mbanga, actuel président de l'orchestre Mokamo Kisina de Ya Mbanga, et de l'accompagner dans sa présence digitale, avec à disposition des réseaux sociaux suivis par plus de 3000 personnes cumulées, une meilleure visibilité lui permettant de rayonner de nouveau dans l'écosystème culturel.

    \begin{figure}[h]
        \centering
        \includegraphics[width=0.9\textwidth]{./media/image2.png}
        \caption{Illustration du clip \og Senda Yima\fg{} de Mokamo Kisina de Ya Mbanga \cite{mbanga_mumbungu_senda_2024}, ayant cumulé plus de 8000 vues sans sponsorisation.}
    \end{figure}

Cette initiative contribue ainsi à corriger les asymétries de
reconnaissance dans les transferts culturelles, en soutenant une justice
épistémique dans la circulation des patrimoines musicaux africains.
\section{Conclusion}\label{conclusion}

L'art musical mbala occupe une place centrale dans la formation de la musique congolaise moderne, mais cette influence demeure largement invisibilisée. Cette invisibilisation relève d'une injustice épistémique structurelle, héritée de cadres théoriques qui ont hiérarchisé les formes de rationalité et disqualifié les traditions orales. Les humanités numériques offrent aujourd'hui une opportunité inédite pour documenter, contextualiser et valoriser ces savoirs, à condition d'être mobilisés dans une perspective critique et inclusive. Ce travail constitue une première étape vers une patrimonialisation numérique du savoir mbala. Il appelle des recherches futures sur les transferts culturels, les cadres juridiques de protection des savoirs locaux, les infrastructures numériques adaptées aux réalités africaines et les méthodologies interdisciplinaires capables d'articuler analyse musicale, anthropologie, philosophie et sciences des données.
%\nocite{*}
\nocite{klein2020webdocumentaire}

\printbibliography

\end{document}